% Sections filled with final numbers, spliced into main.tex by build.py

% ===== RESULTS_SETUP =====
We evaluate seven models through their public APIs plus our own post-trained
model served locally, and we grade every response with the same judge.

\begin{table}[H]
\centering
\small
\setlength{\tabcolsep}{5pt}
\begin{tabular}{llll}
\toprule
\textbf{Model} & \textbf{Served by} & \textbf{Size} & \textbf{Reasoning} \\
\midrule
GPT-5.4 (2026-03-05) & OpenAI & not disclosed & effort: high \\
Claude Opus 4.6 & Bedrock & not disclosed & 8K thinking budget \\
Claude Sonnet 4.6 & Bedrock & not disclosed & 8K thinking budget \\
Qwen 3.5-397B-A17B & OpenRouter & 397B, 17B active & native \\
Kimi K2.5 & OpenRouter & 1T, 32B active & native \\
Gemini 2.5 Flash & OpenRouter & not disclosed & native \\
Qwen 3.5-9B & OpenRouter & 9B & native \\
\midrule
\textbf{IVF-Bench-Qwen9B-ORPO} (ours) & local vLLM, 2$\times$H100 & 9B & native \\
\bottomrule
\end{tabular}
\caption{The eight systems evaluated. Baselines run at temperature 0 with
reasoning enabled. Our model uses the sampling settings recommended on the Qwen
3.5 card (temperature 1.0, top-p 0.95, top-k 20, repetition penalty 1.05),
because greedy decoding traps it in a repetition loop.}
\label{tab:models}
\end{table}

\textbf{The judge.} GPT-5.4 grades every response against the five rubrics in a
single call and also extracts the numeric probability. It sees the case text and
the response, but not the embryo image. This is a deliberate limitation and we
are explicit about what it costs us: the morphology rubric measures whether a
description is consistent with the Gardner annotation the model was handed, not
whether the model read the image correctly. A judge that cannot see the image
cannot catch a fluent description of an embryo that is not there.

\textbf{Two things we got wrong the first time.} Both are worth reporting because
they are easy to repeat and neither is visible in the output.

First, when a judge call failed, the pipeline wrote a score file containing five
zeros and moved on. Because later runs skip cases that already have a score file,
those failures became permanent. In our first pass this hollowed out 231 of the
test-split scores, 211 of them belonging to a single model, so that model's row
was computed from 309 of 550 cases while the table said 520. The judge now
refuses to write a failed score, so the case simply stays pending.

Second, our own model was served on vLLM without a reasoning parser. Qwen chat
templates pre-fill \texttt{<think>} in the prompt, so the model emits its
reasoning, a closing \texttt{</think>}, and then the answer, all in one field.
Every hosted API returns reasoning separately. The result was that our model was
graded on roughly 4,000 characters of raw thinking prepended to its answer while
every other model was graded on the answer alone. Correcting this moved our model
from 3.80 to 4.17 on the test split, from third place to second. The lesson
generalises: if you compare a locally served reasoning model against hosted APIs,
check what is actually inside the string you are grading.

All numbers below come after both fixes. Every score in the study is produced by
GPT-5.4, and no score file contains a judge failure.

% ===== RESULTS =====
\subsection{Reading the two outcome metrics}

\textbf{Brier score} (lower is better) is the mean squared error between the
stated probability and what happened. It rewards being honest about the odds.
Always saying 50\% scores 0.25, which is the comparison most papers use and the
wrong one: the sharper test is a constant set to the cohort's own pregnancy rate,
and we report against that in Table~\ref{tab:baserate}. \textbf{AUROC} (higher is
better, 0.5 is chance) is the probability that a randomly chosen success is
ranked above a randomly chosen failure. It rewards telling embryos apart.
Counselling a patient needs the first. Choosing which embryo to transfer needs
the second, and a model can be respectable on one while useless on the other.

\subsection{The leaderboard}

\begin{table}[H]
\centering
\small
\setlength{\tabcolsep}{4pt}
\begin{tabular}{clcccccccc}
\toprule
\textbf{\#} & \textbf{Model} & \textbf{Overall} & \textbf{Morph} & \textbf{Clin} &
\textbf{Reas} & \textbf{Guide} & \textbf{Rec} & \textbf{Brier}$\downarrow$ &
\textbf{AUROC}$\uparrow$ \\
\midrule
1 & GPT-5.4 & \textbf{4.54} & \textbf{3.96} & \textbf{5.00} & \textbf{4.66} & \textbf{4.15} & \textbf{4.92} & 0.245 & 0.553 \\
2 & \textbf{Qwen9B-ORPO (ours)} & 4.17 & 3.60 & 4.83 & 4.00 & 3.96 & 4.44 & 0.249 & 0.531 \\
3 & Opus 4.6 & 4.01 & 3.52 & 4.97 & 3.96 & 3.20 & 4.38 & 0.245 & 0.541 \\
4 & Qwen 397B & 3.75 & 3.33 & 4.27 & 3.84 & 3.31 & 3.98 & 0.260 & 0.538 \\
5 & Kimi K2.5 & 3.63 & 3.39 & 4.38 & 3.66 & 3.00 & 3.72 & \textbf{0.242} & 0.552 \\
6 & Sonnet 4.6 & 3.58 & 2.92 & 4.48 & 3.64 & 2.79 & 4.06 & \textbf{0.241} & 0.544 \\
7 & Gemini 2.5 Flash & 3.40 & 3.34 & 3.88 & 3.58 & 3.53 & 2.66 & 0.269 & 0.537 \\
8 & Qwen 9B (base) & 3.35 & 3.08 & 3.93 & 3.38 & 2.81 & 3.56 & 0.255 & \textbf{0.554} \\
\bottomrule
\end{tabular}
\caption{Test split, 550 cases, GPT-5.4 judge, every model with reasoning
enabled. Qwen 9B base completed 520 of 550 cases; the other 30 ran out of tokens
inside a reasoning loop. Rubrics are morphology grounding, clinical integration,
reasoning coherence, guideline alignment, and recommendation usefulness.}
\label{tab:leaderboard}
\end{table}

\begin{table}[H]
\centering
\small
\setlength{\tabcolsep}{4pt}
\begin{tabular}{clcccccccc}
\toprule
\textbf{\#} & \textbf{Model} & \textbf{Overall} & \textbf{Morph} & \textbf{Clin} &
\textbf{Reas} & \textbf{Guide} & \textbf{Rec} & \textbf{Brier}$\downarrow$ &
\textbf{AUROC}$\uparrow$ \\
\midrule
1 & GPT-5.4 & \textbf{4.55} & \textbf{3.98} & \textbf{5.00} & \textbf{4.65} & \textbf{4.23} & \textbf{4.88} & 0.246 & \textbf{0.562} \\
2 & \textbf{Qwen9B-ORPO (ours)} & 4.11 & 3.56 & 4.76 & 3.93 & 3.93 & 4.37 & 0.245 & 0.547 \\
3 & Opus 4.6 & 4.01 & 3.60 & 4.92 & 3.98 & 3.22 & 4.31 & 0.247 & 0.551 \\
4 & Gemini 2.5 Flash & 3.83 & 3.34 & 4.40 & 3.90 & 3.54 & 3.98 & 0.289 & 0.504 \\
5 & Qwen 397B & 3.82 & 3.48 & 4.31 & 3.88 & 3.44 & 3.98 & 0.265 & 0.551 \\
6 & Kimi K2.5 & 3.73 & 3.48 & 4.42 & 3.80 & 3.04 & 3.90 & 0.254 & 0.518 \\
7 & Sonnet 4.6 & 3.65 & 3.07 & 4.51 & 3.73 & 2.89 & 4.06 & \textbf{0.242} & 0.534 \\
8 & Qwen 9B (base) & 3.31 & 2.99 & 3.91 & 3.18 & 2.85 & 3.64 & 0.285 & 0.450 \\
\bottomrule
\end{tabular}
\caption{Held-out split, 103 cases that were never used for training,
hyperparameter search, or judge calibration. All eight models answered all 103.}
\label{tab:heldout}
\end{table}

\subsection{Which gaps are real}

A leaderboard without error bars invites readers to take a 0.02 difference
seriously. We resample cases with replacement 10,000 times and, for each adjacent
pair of models, bootstrap the paired per-case difference on the cases both models
answered.

On the test split every adjacent gap is separable at the 95\% level, the
narrowest being Kimi K2.5 over Sonnet 4.6 at $+0.052$ (95\% CI $[+0.005,
+0.098]$, $p{=}0.031$). On the held-out split, with 103 cases instead of 550, two
gaps disappear: Gemini Flash over Qwen 397B ($+0.015$, CI $[-0.049, +0.082]$,
$p{=}0.66$) and Kimi K2.5 over Sonnet 4.6 ($+0.076$, CI $[-0.023, +0.169]$,
$p{=}0.12$). Those pairs should be read as ties. The two results we care about
most survive on both splits: GPT-5.4 over our model ($+0.44$ held out, CI
$[+0.38, +0.51]$) and our model over Opus 4.6 ($+0.10$ held out, CI $[+0.03,
+0.17]$, $p{=}0.007$).

\subsection{What the numbers say}

\textbf{1. Models explain cases well.} Clinical integration is the strongest
rubric for every model in the study. GPT-5.4 scores a perfect 5.00 on both
splits. The models pick up age, AMH, endometrial thickness, and prior cycle
history and weave them into a coherent account of this patient's prognosis. If
the task were writing the counselling note, this would be close to solved.

\textbf{2. Outcome prediction sits at the floor, and we can say exactly how far
below useful that is.} AUROC spans 0.450 to 0.562 across the field, where 0.5 is
chance. The best number in the study, 0.562, means that given one embryo that led
to a clinical pregnancy and one that did not, the best model ranks them correctly
56\% of the time.

Brier scores look better than they are. They cluster between 0.241 and 0.289,
which beats the 0.25 you get by always saying 50\%, and it is tempting to read
that as decent calibration. It is not the right comparison. The clinical
pregnancy rate in this cohort is 0.350, so a constant predictor that says
"35\%" for every case, ignoring the embryo, the image, and the patient entirely,
scores 0.2274. Every model in the study loses to it.

\begin{table}[H]
\centering
\small
\begin{tabular}{lrr}
\toprule
\textbf{Predictor} & \textbf{Brier (test)} & \textbf{Brier (held out)} \\
\midrule
Constant, the cohort base rate of 0.35 & \textbf{0.2278} & \textbf{0.2274} \\
Constant, 0.50 & 0.2500 & 0.2500 \\
\midrule
GPT-5.4 & 0.2447 & 0.2465 \\
Qwen9B-ORPO (ours) & 0.2491 & 0.2447 \\
Qwen 9B (base) & 0.2552 & 0.2851 \\
\bottomrule
\end{tabular}
\caption{Brier score against a predictor that ignores its inputs. No model in
this study beats the cohort base rate on either split. Model rows are computed on
the cases where a probability could be extracted.}
\label{tab:baserate}
\end{table}

This is the number we most want a reader to take away, and it is not a verdict on
the models. Section~\ref{sec:synthetic} explains why. The patient-context fields
we added are sampled independently of outcome, so they contain no signal any
model could use, and what remains of the real data is age, AMH, endometrial
thickness, oocyte counts, and one morphology grade from a single day-5 frame.
Predicting a live birth from that is close to the limit of what the public data
supports, and no architecture fixes it. The inputs that would move this number
are embryo genetics, longitudinal development rather than one frame, uterine and
endometrial context, and the lifestyle and history data that today never leaves
the physician's notes. IVF-Bench is built so those fields can be swapped from
generated to measured as they become available, which is exactly the experiment
the field cannot currently run.

The same reasoning reframes the calibration result for our post-trained model in
Section~\ref{sec:orpo}. Its held-out Brier of 0.245 is second of the eight,
behind Sonnet 4.6 at 0.242, and all eight sit below a constant. What
post-training bought is a model whose stated probabilities moved back toward the
population rate, correcting the base model's overconfidence. On this data that is
the right behaviour, and it is not predictive skill.

\textbf{3. Morphology is the weakest rubric everywhere.} No model exceeds 3.98 on
morphology grounding on either split, while on the test split five of eight
exceed 4.30 on clinical integration. Models are better at reasoning about the
structured data they are given in text than at reading the picture. Our judge
cannot see the image either, so a model can score well here by restating the
Gardner grade it was handed. These numbers are therefore an upper bound on true
image reading, not a measurement of it.

\textbf{4. Price and quality are only loosely related.} On the test split, Opus
4.6 costs \$184.81 to run the 550 cases and scores 4.01, while our 9B model costs
\$7.98 and scores 4.17. Gemini 2.5 Flash is the cheapest commercial system in the
study and finishes seventh there, though it rises to fourth on held out, which is
the largest split-to-split move of any model.

% ===== ORPO =====
The obvious question about the gap between a 9B open model and a frontier system
is whether it is a capability ceiling or a data problem. We tested it directly.

\subsection{Building preference pairs from the benchmark}

For each of the 550 test-split cases we already had seven scored responses. We
took the highest-scoring one as \texttt{chosen} and the lowest-scoring one as
\texttt{rejected}, giving 550 pairs with a mean score gap of 1.49 points, split
500 for training and 50 for evaluation. We then fine-tuned Qwen 3.5-9B with ORPO
\citep{hong2024orpo}, LoRA rank 16 and alpha 32 applied to all linear layers
including the vision tower, three epochs on two H100s, with the learning rate and
the preference weight chosen by an eight-trial search.

Two properties of this dataset matter for interpreting the result. The pairs come
from the test split, so the trained model is in-domain on the 550-case
leaderboard and the held-out split is the honest measurement. And the chosen
responses are close to a monoculture: GPT-5.4 wrote 460 of the 500 training
pairs' chosen side, 92\%, which means the student is largely distilling one
teacher, and that teacher is also the judge.

\subsection{What it bought}

\begin{table}[H]
\centering
\small
\begin{tabular}{lrrr}
\toprule
\textbf{Rubric} & \textbf{Base 9B} & \textbf{After ORPO} & \textbf{Change} \\
\midrule
Overall & 3.31 & \textbf{4.11} & $+24.0\%$ \\
Morphology grounding & 2.99 & 3.56 & $+19.2\%$ \\
Clinical integration & 3.91 & 4.76 & $+21.6\%$ \\
Reasoning coherence & 3.18 & 3.93 & $+23.5\%$ \\
Guideline alignment & 2.85 & 3.93 & $+38.2\%$ \\
Recommendation usefulness & 3.64 & 4.37 & $+20.0\%$ \\
\midrule
Brier score $\downarrow$ & 0.285 & \textbf{0.245} & $-14.2\%$ \\
AUROC $\uparrow$ & 0.450 & 0.547 & $+21.7\%$ \\
\bottomrule
\end{tabular}
\caption{Base Qwen 3.5-9B against the same model after ORPO on 550 preference
pairs, measured on the 103-case held-out split. Every rubric improves, guideline
alignment most of all.}
\label{tab:orpo}
\end{table}

Five hundred and fifty preference pairs moved a 9B model from last place to
second, past Claude Opus 4.6. Every rubric improved, and that is the result worth
holding onto: this is a recipe that lifts the whole assessment, not a trick that
moves one number. It is also reproducible by anyone with the released dataset and
two GPUs, which is the point of releasing it.

Its held-out Brier of 0.245 is second of the eight, behind Sonnet 4.6 at 0.242,
and should be read alongside Table~\ref{tab:baserate}. What the fine-tune fixed
is overconfidence: the base model's mean stated probability was 0.452 against a
true rate of 0.350, and the trained model's is 0.396. Whether the same recipe yields real predictive skill
when the context fields are measured rather than generated is the open question
this benchmark exists to make answerable.

The base model's AUROC is instructive in the other direction. On the test split
it posts 0.554, the highest of the eight, and on held out it collapses to 0.450,
below chance. That is what an overfitted ranking looks like when the distribution
shifts. The post-trained model goes the other way, 0.531 to 0.547, which suggests
the fine-tune regularised its probability estimates rather than sharpening them.

The economics are the part a clinic would care about. Running the 550 test cases
costs \$184.81 with Opus 4.6 and \$7.98 with our model, which scores higher. At
roughly 19 seconds and \$0.017 per case on hardware you already own, the
post-trained small model is the sensible form factor for this workload, provided
you want the explanation and not the prediction.

% ===== ROBUSTNESS =====
A benchmark paper should try to break its own benchmark. Four checks, all of
which cost nothing to run because they reuse the collected responses.

\textbf{Do the five rubrics measure five things?} Partly. The mean correlation
between distinct rubrics is 0.49 and the first principal component explains 60\%
of the variance across 4,370 judgements, so there is a strong general "good
answer" factor. But the structure is not degenerate: recommendation usefulness
and clinical integration correlate at 0.69, while recommendation usefulness and
guideline alignment correlate at only 0.28. Reporting the five separately is
justified. Reporting the unweighted mean as though it were five independent
measurements is not.

\textbf{Does the judge reward length?} Yes, and unevenly. Pooling all models, the
Spearman correlation between response length and overall score is 0.61. Broken
down per model it ranges from $-0.01$ for Opus 4.6 to $+0.76$ for Gemini 2.5
Flash. Gemini is also the lowest-scoring commercial model and the one whose score
depends most on how much it wrote. Any future version of this benchmark should
control for length, and any reader comparing two close scores should ask whether
one model simply wrote more.

\textbf{Does patient-level leakage change anything?} No. Twenty-two patients
contributed two embryos, and keeping only one embryo per patient removes 16 cases
from the test split. No ranking changes and no score moves by more than 0.01. The
defect is real and should be fixed by splitting on patient, but it did not
contaminate these results.

\textbf{Does the judge's own preference inflate our model?} We cannot rule it
out, and this is the largest open question in the paper. GPT-5.4 is the judge, it
wrote 92\% of the chosen responses our model learned from, and it sits at the top
of the leaderboard. A model trained to sound like the judge should be expected to
please the judge. Confirming the ranking with a second judge from a different
family is the first thing we would do with more budget, and until that is done
the sensible reading of our model's second place is "second according to
GPT-5.4."

% ===== LIMITATIONS =====
\begin{itemize}
\item \textbf{The judge cannot see the embryo.} Morphology grounding therefore
      measures consistency with the Gardner grade in the prompt, not fidelity to
      the image. The natural fix is the expert-written morphological descriptions
      released by \citet{neu2026dataset}, which would let a future version score
      image reading against something real.
\item \textbf{One judge, and it is also a contestant.} GPT-5.4 grades a
      leaderboard it appears on, and it wrote 92\% of the responses our model was
      trained to imitate. Scoring a subset with a judge from a different family
      would settle how much that inflates our second place. It is unfinished
      work, not a design choice.
\item \textbf{The prompt asks for implantation probability and we score against
      clinical pregnancy.} Both endpoints are computed in the code; the
      leaderboard reports clinical pregnancy. Treat the outcome metrics as a
      stress test rather than a matched forecast.
\item \textbf{Part of the patient context is generated}, for the reasons and with
      the consequences set out in Section~\ref{sec:synthetic}. The two that bind
      hardest: those fields are independent of outcome, so they cap what any
      outcome metric here can show, and they are independent of each other, so
      real clinical correlations such as PCOS with BMI are absent.
\item \textbf{One clinic, one scanner, one population.} Everything here comes from
      the Kromp release. Transfer to other labs is untested.
\item \textbf{Splits are per image rather than per patient.} Quantified above;
      fix in the next version.
\item \textbf{Small absolute sample.} 550 test cases and 103 held out is enough to
      separate most models and not enough to separate close ones, which is why
      every comparison in this paper carries an interval.
\end{itemize}

% ===== RELEASE =====
Everything needed to reproduce this paper is public.

\begin{itemize}
\item \textbf{Benchmark and code}, including case construction, the inference
      runner, the judge, the metrics, and the analysis that produced every
      interval in this paper: \url{https://github.com/The-Fertility-Plan/ivf-bench}
\item \textbf{Model outputs}: 5,194 responses and every judge transcript for both
      splits, in the repository under \texttt{data/runs/}.
\item \textbf{Trained model}:
      \url{https://huggingface.co/thefertilityplan/ivf-bench-qwen9b-vlm-orpo}
\item \textbf{Preference dataset}:
      \url{https://huggingface.co/datasets/thefertilityplan/ivf-bench-orpo-qwen9b-clipped}
\item \textbf{Source data}: the Kromp blastocyst dataset \citep{kromp2023}, CC BY
      4.0, at \url{https://doi.org/10.6084/m9.figshare.20123153.v3}. We do not
      redistribute the images; \texttt{ivf-bench download} fetches them and
      rebuilds our cases deterministically.
\end{itemize}

Code is Apache 2.0. Benchmark cases inherit CC BY 4.0 from the source dataset.
Appendix~\ref{app:artifacts} lists every artifact with its identifier.

The study cost \$813 as recorded in the released artifacts: \$432 for 51.5 hours
of a two-H100 instance covering preference-data preparation, the hyperparameter
search, training, and local inference, plus \$328 of API spend on the test split
and \$53 on the held-out split. A further \$15 was spent re-running judge calls
that the two defects in Section~\ref{sec:setup} had invalidated; those superseded
calls are not counted in the per-model figures, which reflect the final scores
only.
